\section{Discussion}\label{planta:sec:discussion}

This section discusses the SUAVE and navigation results, explains the choice of OWL rather than PDDL for modeling the system architecture, and examines the limitations of \acronym.

\subsection{SUAVE exemplar}

The SUAVE experiments show that applying \acronym to a ROS~2-based robotic exemplar is feasible, but they also reveal mission-performance limitations relative to the other managing subsystems evaluated. BT achieved lower \emph{pipeline search times} and greater \emph{pipeline distances inspected} than \acronym across the original and extended experiments, and ROSA and Metacontrol also outperformed \acronym on several mission-performance comparisons. Therefore, the SUAVE results should be interpreted primarily as evidence that \acronym can be integrated with and executed on a realistic ROS~2-based robotic exemplar, rather than as evidence of mission-performance superiority. The SUAVE exemplar provides a concrete scenario and an experimental setup for evaluating whether managing systems can bridge the gap between theory and practice. The challenges involved in achieving this integration are often underestimated. As shown in~\Cref{planta:sec:impl_plansys}, integrating \acronym with \plansys required extending the PDDL features that \plansys supports, adapting its plan-execution algorithms, and implementing a planner plugin, which introduced substantial engineering overhead.

At the same time, the SUAVE experiment highlights the need for new ROS~2-based exemplars for self-adaptive robotic systems that enable a more fine-grained evaluation of how overall mission quality is affected by performance efficiency, including scalability, reaction time, and memory usage.

\subsection{Navigation use case}\label{planta:sec:disc_nav}
Although both \acronym and the approach proposed by Cámara et al.~\cite{CAMARASoftwareArchitectureTask2020} address TACA for the entire mission, they differ fundamentally in their underlying planning philosophy. \acronym computes a satisficing task and configuration plan that selects the best configuration for each action based on a predefined utility function. In contrast, the approach by Cámara et al.\ seeks to generate an optimal task and configuration policy that maximizes the overall mission utility. Furthermore, \acronym assumes deterministic action effects and configuration utilities, while the method of Cámara et al.\ models probabilistic outcomes, such as collision risks during execution. % chktex 12 -- sentence starts with a macro

In the experiment in \Cref{planta:sec:nav_exp}, the replicated final step of \CamaraApproach, which computes an optimal policy under a probabilistic model, required more planning time than \acronym. Optimal planning can yield better results when the robot and its environment behave as modeled, but in dynamic environments the assumptions made at design time may not hold during operation. Under these conditions, it can be more advantageous to search for a satisficing solution, monitor during execution whether the current plan remains valid, and replan when it does not, because optimal planning is generally more expensive and therefore incurs a higher penalty at each replanning. A practical strategy could therefore be to use optimal planning for the initial mission plan and satisficing planning for runtime replanning.

\subsection{Modeling the System Architecture in OWL}
There are three main reasons why the robotic system and its functional designs are modeled in OWL rather than directly in PDDL\@.

First, although PDDL supports type hierarchies, OWL offers substantially richer mechanisms for representing class axioms, relations, and inferred knowledge. In PDDL, objects are assigned a type, and type membership follows from explicitly declared type hierarchies. This limits the representation of relationships between classes, which is particularly relevant for robotic components. For example, a specific camera model can be represented as a camera, a visual sensor, and a device requiring electrical power, but these relationships must be explicitly encoded in the type hierarchy. Moreover, it remains unclear whether PDDL types can be treated as predicates~\cite{GOLDMANpddlTypeProblems2012}. In contrast, OWL allows classes to be defined independently and related through axioms, enabling class membership to be inferred rather than explicitly declared. For example, a device requiring electrical power can be defined as any component with a power-consumption property, allowing camera models with this property to be classified accordingly. Together with SWRL rules, OWL also supports inferring relations that are not explicitly asserted. \acronym relies on these inference capabilities to determine which $\FunctionDesigns$ can ground a $\FunctionTomasys$ and which $\FunctionDesign$ has higher utility, as described in \Cref{planta:sec:tomasys_pddl}. These modeling capabilities are particularly valuable for representing robotic systems and are extensively used in \acronym. % chktex 12 -- sentence starts with a macro

Secondly, there exist many tools for OWL ontologies that assist in modeling classes and their relationships, including the SWRL rules used in this work for inference. These tools allow automatic checking of the model for common errors, such as verifying the consistency of class descriptions and ensuring that instances of the classes can exist without violating other axioms (i.e., satisfiability). In addition, they can visualize the class structure, which further facilitates model development.

Lastly, established standards exist for representing robotic systems and their components using ontologies, such as the IEEE Standard for Autonomous Robotics (AuR) ontology~\cite{IEEEStandardAutonomous2022}. No equivalent standards are known to exist for PDDL\@. Therefore, modeling with OWL ontologies increases the likelihood of producing reusable and interoperable models.

\subsection{Limitations}\label{planta:sec:limitations}
A limitation of \acronym is that it evaluates architectural alternatives locally, i.e., for each $\FunctionTomasys$ required by an action, it selects the realizable $\FunctionDesign$ with the highest utility with respect to that $\FunctionTomasys$. This selection should not be interpreted as mission-level utility optimization. For example, an architectural configuration that provides better performance for an individual function may consume significantly more energy throughout the mission, yet \acronym would still select it.
In domains where architectural choices have long-term effects, such as energy consumption or resource depletion, accounting for such trade-offs would require explicitly modeling their mission-level consequences and incorporating an appropriate optimization objective.

A further limitation concerns the costs associated with architectural reconfiguration. Although reconfiguration is represented through planning actions, the current formulation does not model transition-related costs, such as reconfiguration duration and energy consumption. Consequently, the resulting plans may overlook relevant overhead associated with changing architectural configurations. These aspects could be incorporated by extending the planning model with temporal or cost-aware features, such as durative actions or action costs.

Another limitation of \acronym is its reliance on PDDL derived predicates, a feature not supported by most PDDL-based planners, limiting the number of planners that can be used with \acronym. More specifically, to the best of our knowledge, no existing planner supports both derived predicates and numeric fluents. This poses a significant constraint in the robotics domain, where numeric modeling is essential for representing properties such as battery level, power consumption, and distance.

To mitigate this limitation, the proposed \owltopddl transformation encodes numerical values as symbolic constants and explicitly represents inequality relations among them. This enables symbolic reasoning over ``numeric'' values using inequalities. However, it does not support traditional numerical planning as defined in PDDL 2.1~\cite{fox2003pddl2}, which involves arithmetic operations and numeric fluents.

This work nevertheless keeps derived predicates explicit in the planning domains of \acronym, because compiling them away can, in the worst case, cause an exponential blow-up in the size of the domain formulation or in the length of the shortest plans~\cite{THIEBAUX200538}.

%%% Local Variables:
%%% mode: latex
%%% TeX-master: "main"
%%% End:
